\phantomsection
\section*{Abstract}
\addcontentsline{toc}{section}{Abstract}

Reading is central to everyday life. However, it can also present new significant challenges for individuals with reading difficulties, such as those with age-related vision loss or dyslexia. For these readers, written text can impose demands that exceed what they can comfortably process, making reading effortless and less accessible.
Digital text can help: it can be reshaped as it is read, and an eye tracker can show when a reader is struggling, so the text can adapt in real time without breaking the reader's flow. This is what the reading the reader (RtR) project is exploring to achieve.
Earlier prototypes in the RtR project showed that achieving the aforementioned goals is definitely possible, but each of them mixed sensing, decision-making, and the reading screen into one tangled codebase, so every new study meant rewriting core code, and none gave researchers a controlled and reproducible approach to run such studies.

This thesis builds and evaluates a modular, researcher-operated framework for adaptive reading.
The participant reads on one screen, where the text adapts to their gaze, while the researcher watches the live gaze points projected on a second screen and can trigger or approve changes as the session runs.
Inside the , four parts (sensing, analysis, decision, and intervention) are separate, swappable modules, so a new intervention or decision-maker, including an AI one, can be added without changing the rest.
When the text changes, the platform keeps the reader's place, and it records every gaze sample, decision, and change so a session can be replayed and reproduced.

We evaluated the platform on a real Tobii eye tracker over eight reading sessions.
New modules and three outside decision providers connected without changing the core; gaze streamed at the rated \SI{90}{\hertz} with validity above \SI{92}{\percent}; and the system met its \SI{100}{\milli\second} latency budget on every sample.
With place-keeping on, readers tended to return to their place faster after a change.
Two limitations remain: the automatic decision path was tested in a single run, not yet at full scale, and a revised place-keeping method has so far been checked only by geometry, not with readers, so its benefit to reading is unconfirmed.
Even so, the platform shows that the adaptive reading loop can be cleanly separated under real-time conditions, and it gives future studies of reading interventions and AI-driven adaptation a solid base to build on.
